\begin{figure}[h!]
    \centering
    \fbox{
    \begin{minipage}{0.95\textwidth}
    \centering
    \small
    \setlength{\jot}{2pt}
    \begin{aligned}
        \text{CSV (location mentions)}
        &\rightarrow \text{Text normalization} \\
        &\rightarrow \text{Geocoding (Nominatim / OpenStreetMap)} \\
        &\rightarrow \text{Latitude \& Longitude + place metadata} \\
        &\rightarrow \text{Filter: Netherlands only + remove country-level results} \\
        &\rightarrow \text{Spatial matching: NUTS3 containment (preferred)} \\
        &\rightarrow \text{Fallback: NUTS2 containment or nearest region (distance-based)} \\
        &\rightarrow \text{Optional expansion: NUTS2 } \rightarrow \text{ NUTS3 mapping (weights)} \\
        &\rightarrow \text{Final enriched dataset (CSV)}
    \end{aligned}
    \end{minipage}
    }
    \caption{Geocoding and spatial enrichment pipeline mapping location mentions to Dutch NUTS regions using hierarchical and distance-based matching.}
    \label{fig:geocoding_pipeline}
\end{figure}
The geocoding pipeline converts location mentions from a CSV dataset into structured geographic data that can be used for spatial analysis of drought impacts. It uses the Nominatim API from OpenStreetMap to translate text-based location names into latitude and longitude coordinates. These coordinates are then linked to Dutch administrative regions (NUTS2 and NUTS3) to enable consistent regional analysis. The overall workflow is summarised in Figure~\ref{fig:geocoding_pipeline}.

Each location mention is first cleaned by normalising whitespace. Duplicate mentions are cached to avoid repeated API requests. A fixed delay between requests is included to comply with API usage limits. The system mainly focuses on locations within the Netherlands by filtering results using country codes.

When multiple geocoding results are returned, the pipeline selects the most relevant one. Preference is given to broader administrative levels such as countries or regions when the input is general. Each result is then classified into a place type such as country, region, city, or local area based on OpenStreetMap metadata.

After geocoding, each point is matched to Dutch NUTS regions. The pipeline first checks whether a point falls within a NUTS3 boundary, which is the preferred level of resolution, or otherwise falls back to NUTS2 boundaries. If no direct match is found, a distance-based method in a projected coordinate system identifies the nearest region within a defined threshold. This ensures that most locations can still be assigned to a meaningful administrative unit.

To improve performance, simplified NUTS boundary files are precomputed and stored locally. A mapping between NUTS2 and NUTS3 regions is also created in advance to support efficient aggregation of results. In cases where only NUTS2 information is available, values can be distributed across related NUTS3 regions using uniform weights.

The final output combines the original dataset with geocoding results and regional assignments. Entries outside the Netherlands or only resolved at the country level are removed to maintain consistent spatial resolution. This design balances accuracy, speed, and robustness, making the dataset suitable for further analysis of drought-related impacts at regional scale.